% 本文件由 LaTeX 排版 Agent 根据有效 translation.json 自动生成。
% author=Gregory Thaumaturgus; work_id=0607; title=论信德的十二个主题

\chapter{论信德的十二个主题}

% structure: p0001 source=h1 target=omitted reason=page-title-already-rendered-by-parent

其中亦给出绝罚公式，并对每条加以阐释。\par

% structure: p0003 source=h2 target=section reason=relative-heading-rank; base=h2

\section{主题一}

若有人说基督的身体是非受造的，并拒绝承认祂——天主的非受造圣言（天主）——取了受造人性的血肉，并成为可见的降生成人者，正如经上所载，应受绝罚。\par

\Emph{阐释}\par

怎么能说身体是非受造的呢？因为非受造的是无情感、不受伤、不可触的。但基督从死者中复活后，向祂的门徒显示了钉痕和枪伤，以及一个可触摸的身体，尽管祂也是在门户关闭时进到他们中间的，为同时向他们显示天主性的能力和身体的真实性。\par

然而，祂虽是天主，却以自然的方式被认作人；祂真实地作为人存在，却也因祂的作为被显示为天主。\par

% structure: p0008 source=h2 target=section reason=relative-heading-rank; base=h2

\section{主题二}

若有人说基督的血肉与天主性是同质的，并拒绝承认祂——在万世之前以天主的形体存在——空虚了自己，取了奴仆的形体，正如经上所载，应受绝罚。\par

\Emph{阐释}\par

怎么能说受时间限制的血肉与永恒的天主性同质呢？因为所谓同质，是指在本质上和永恒延续上相同，毫无变易。\par

% structure: p0012 source=h2 target=section reason=relative-heading-rank; base=h2

\section{主题三}

若有人说基督如同先知中的一位，取了完满的人，并拒绝承认祂——因血肉由童贞女所生——成了人，生于白冷，在纳匝肋长大，年龄增长，满了规定的年数（公开出现）并在约旦河受洗，且从圣父领受这见证：\Quote{这是我的爱子，}\BibleRef{玛 3:17}，正如经上所载，应受绝罚。\par

\Emph{阐释}\par

怎么能说基督（主）如同先知中的一位取了完满的人呢？因为祂本身就是主，藉由童贞女实现的降生成人而成了人。因此经上记载：\Quote{第一个人是出于地，属于土。}\BibleRef{格前 15:47} 但出于土的归于土，而成为第二个人的归于天。我们也读到\Quote{第一个亚当和最后一个亚当。}\BibleRef{格前 15:45} 既然第二个依肉体来自第一个，为此基督也被称为人和人子；同样也证实第二个是第一个的救主，祂为此从天降下。既然圣言从天降下，成了人，又升到天上，因此祂被称为第二个来自天上的亚当。\par

% structure: p0016 source=h2 target=section reason=relative-heading-rank; base=h2

\section{主题四}

若有人说基督是由童贞女从人的种子而生，如同所有人一样，并拒绝承认祂是由圣神和童贞玛利亚成了血肉，并由达味的后裔成为人，正如经上所载，应受绝罚。\par

\Emph{阐释}\par

怎么能说基督是由童贞女从人的种子而生呢？因为神圣福音和天使在传报喜讯时，为童贞玛利亚作证说：\Quote{这事如何成就？因为我不认识男人。}\BibleRef{路 1:34} 因此他说：\Quote{圣神要临于你，至高者的能力要庇荫你，因此那要诞生的圣者，将称为至高者的儿子。}\BibleRef{路 1:35} 又对若瑟说：\Quote{不要怕娶你的妻子玛利亚，因为那在她内受孕的是出于圣神。她要生一个儿子，你要给祂起名叫耶稣，因为祂要把自己的民族从他们的罪恶中拯救出来。}\BibleRef{玛 1:20}\par

% structure: p0020 source=h2 target=section reason=relative-heading-rank; base=h2

\section{主题五}

若有人说在万世之前的天主子与在这些末期显现的那位是不同的，并拒绝承认在万世之前的那位与在这些末期显现的那位是同一的，正如经上所载，应受绝罚。\par

\Emph{阐释}\par

怎么能说在万世之前的天主子与在这些末期显现的那位是不同的呢？因为主亲自说：\Quote{在亚巴郎出现以前，我就存在。}\BibleRef{若 8:58} 又说：\Quote{我由天主而来，并来到世上；我又去往父那里。}\par

% structure: p0024 source=h2 target=section reason=relative-heading-rank; base=h2

\section{主题六}

若有人说受苦难的那位与不受苦难的那位是不同的，并拒绝承认圣言——祂自身是无情感、不变的天主——在祂所取的血肉中真实地、但毫无改变地受了苦难，正如经上所载，应受绝罚。\par

\Emph{阐释}\par

怎么能说受苦难的那位与不受苦难的那位是不同的呢？因为主亲自说：\Quote{人子必须受许多苦，被杀害，第三天从死者中复活。}\BibleRef{玛 16:21} 又说：\Quote{你们要看见人子坐在父的右边。}\BibleRef{玛 26:64} 又说：\Quote{当人子在他父的荣耀中来临时。}\BibleRef{玛 16:27}\par

% structure: p0028 source=h2 target=section reason=relative-heading-rank; base=h2

\section{主题七}

若有人说基督是被拯救的，并拒绝承认祂是世界的救主和世界的光，正如经上所载，应受绝罚。\par

\Emph{阐释}\par

怎么能说基督是被拯救的呢？因为主亲自说：\Quote{我是生命。}又说：\Quote{我来是为使他们获得生命。}\BibleRef{若 10:10} 又说：\Quote{信我的人，必不见死亡，而必得见永生。}\BibleRef{若 8:51}\par

% structure: p0032 source=h2 target=section reason=relative-heading-rank; base=h2

\section{主题八}

若有人说基督是完满的人，又是天主圣言，但两者是分离的，并拒绝承认一位主耶稣基督，正如经上所载，应受绝罚。\par

\Emph{阐释}\par

怎么能说基督是完美的人又是圣言天主，却又将两者分开呢？主亲自说：\BibleQuote{你们为什么想要杀死我，这个把从天主那里听见的真理告诉你们的人？} \BibleRef{若 8:40} 因为圣言天主并没有为我们献上一个人，而是为我们献上了祂自己，为了我们而成为人。因此祂说：\BibleQuote{你们拆毁这座圣殿，三天之内我要把它重建起来。} 但祂所说的是祂身体的圣殿。\par

% structure: p0036 source=h2 target=section reason=relative-heading-rank; base=h2

\section{第九论}

若有人说基督遭受变化或改易，并拒绝承认祂在圣神内是不变的，虽有血肉的可朽坏性 \BibleRef{若 2:20}，则让他受诅咒。\par

\Emph{说明}\par

怎么能说基督遭受变化或改易呢？主亲自说：\BibleQuote{我是不改变的；祂的灵魂不会被遗弃在阴间，祂的肉身也不会见朽坏。} \BibleRef{拉 3:6}\par

% structure: p0040 source=h2 target=section reason=relative-heading-rank; base=h2

\section{第十论}

若有人断言基督只是局部地取了人，并拒绝承认祂除了罪以外，在一切事上都与我们相似，则让他受诅咒。\par

\Emph{说明}\par

怎么能说基督只是局部地取了人呢？主亲自说：\BibleQuote{我为羊舍弃我的生命，为的是再取回它来} \BibleRef{若 10:17}，又说：\BibleQuote{我的肉是真实的食粮，我的血是真实的饮料} \BibleRef{若 6:55}，又说：\BibleQuote{谁吃我的肉，并喝我的血，就有永生} \BibleRef{若 6:56}\par

% structure: p0044 source=h2 target=section reason=relative-heading-rank; base=h2

\section{第十一论}

若有人断言基督的身体没有灵魂和理智，并拒绝承认祂是完美的人，在一切事上与我们完全相同，则让他受诅咒。\par

\Emph{说明}\par

怎么能说主（基督）的身体没有灵魂和理智呢？因为不安、忧伤和痛苦，既不是没有灵魂的血肉的特性，也不是没有理智的灵魂的特性；它们既不是不变神性的标志，也不是纯粹幻影的标记，也不标志人性的软弱；而是圣言在祂自身内显出了我们本性的情感和感受的运作，祂穿上了我们的可受苦性，正如经上所记载：\BibleQuote{祂背负了我们的忧患，承担了我们的痛苦。} \BibleRef{依 53:4} 因为不安、忧伤和痛苦，是灵魂的紊乱；而劳苦、睡眠和身体易受伤，是血肉的软弱。\par

% structure: p0048 source=h2 target=section reason=relative-heading-rank; base=h2

\section{第十二论}

若有人说基督只是以幻象在世上显现，并拒绝承认祂确实降生在肉身中，则让他受诅咒。\par

\Emph{说明}\par

怎么能说基督只是以幻象在世上显现呢？祂诞生在白冷，受行割礼，被西默盎举起，养到十二岁（在家），顺从父母，在约旦河受洗，被钉在十字架上，又从死者中复活。\par

因此，当说到祂\BibleQuote{心神烦乱}、祂\BibleQuote{心里忧伤} \BibleRef{玛 26:38}、祂\BibleQuote{身体受伤} \BibleRef{依 53:5} 时，祂向我们展示我们本性所固有的易感受的名称，为表明祂在世上成为人，与世人来往 \BibleRef{巴 3:38}，却没有罪。因为祂按血肉诞生在白冷，其方式符合神性：天上的天使认出祂是他们的主，向那时裹在襁褓中躺在马槽里的祂高唱赞歌，称祂为他们的天主，喊道：\BibleQuote{天主受享光荣于高天，良善的人在世享平安。} \BibleRef{路 2:14} 祂在纳匝肋长大；但以神圣的方式，祂坐在经师中间，以超乎年龄的智慧令他们惊讶，这是按祂肉身生活的能力，正如福音叙事所记载的。祂在约旦河受洗，不是为自己接受任何圣化，而是将圣化的参与赐给他人。祂在旷野受试探，不是屈服于试探，而是面对试探者的挑战时，将我们的试探放在自己面前，以显示试探者的无力。\par

因此祂说：\BibleQuote{你们放心，我已战胜了世界。} \BibleRef{若 16:33} 祂这样说，不是向我们展示只属于天主的任何争斗，而是表明我们自己的血肉有能力战胜苦难、死亡和朽坏，为的是：既然罪恶藉着血肉进入了世界，死亡藉着罪恶统治了众人，那么血肉中的罪恶也能藉着同一血肉的相似而被定罪；并且罪恶的监督者——那试探者——可以被战胜，死亡从它的统治中被推翻，身体埋葬中的朽坏被消除，复活的初果被显示，正义的原则藉着信德开始在世上运行，天国被传报给人，人与天主之间建立了共融。\par

为了这恩宠，让我们光荣圣父，祂为了世界的生命赐下了祂的独生子。让我们光荣那在我们内运行、使我们活过来、并赐予适于与天主共融的恩赐的圣神；让我们不要以死寂的辩论来干预福音的圣言，散布无休止的质疑和言辞之争，使温柔而简单的信德之言变得艰难；而让我们实行信德的工作，热爱和平，展示和谐，保持合一，培养爱德，因为天主喜欢爱德。\par

既然父以自己的权柄所定的时候和日期，不是我们应当知道的，\BibleRef{宗 1:7}而只是相信时间必有终结，将来世界必得显现，审判必得启示，天主子必来降临，行为必得报赏，并承受天国为产业，同样，天主子如何成为人，也不是我们应当知道的；因为这是一个伟大的奥秘，正如经上所载：\BibleQuote{谁能述说他的世代？因为他的生命从地上被夺去。}\BibleRef{依 53:8}我们只应相信，天主子按圣经成了人；他按圣经在世上显现，与世人交往，相似他们，却没有罪；他按圣经为我们死了，又从死者中复活；他升了天，坐在圣父的右边，按圣经他要从那里降来，审判生者死者；免得我们彼此以言语争战，使某些人亵渎信仰之言，以致应验经上所载：\BibleQuote{因了你们，我的名在异民中不断受亵渎。}\BibleRef{依 52:5}\par

\SourceRecord{Translated by S.D.F. Salmond.；Ante-Nicene Fathers；Vol. 6.；Edited by Alexander Roberts, James Donaldson, and A. Cleveland Coxe.；Buffalo, NY: Christian Literature Publishing Co.,；1886.；<http://www.newadvent.org/fathers/0607.htm>.}
